%%%%%%%%%%  scan idx 220  |  folio 207  %%%%%%%%%%
\emph{Preuve.} Le fait que $i_s^{-1}\underline{\Homop}(X,Y)$ soit mou résulte de (2.2) appliqué pour $A=\varnothing$; prenant $S=s$, on en déduit que $\underline{\Homop}(X_s,Y_s)$ est mou. La dernière assertion découle de (2.2).

\textbf{Corollaire 2.4.} Sous les hypothèses de (2.2), soit $E$ un $\cO_X$-Module. Pour tout fermé $A$ de $X_s$, l'homomorphisme canonique
\[
  H^0\bigl(X_s,i_s^{-1}(E)\bigr)
  \longrightarrow
  H^0\bigl(A,i_s^*(E)\bigr),
\]
déduit de l'épimorphisme
\[
  i_s^{-1}(E)\longrightarrow i_s^*(E)
  =i_s^{-1}(E)\otimes_{i_s^{-1}(\cO_X)}\cO_{X_s},
\]
est un épimorphisme.

\emph{Preuve.} Comme $i_s^{-1}(\cO_X)$ est mou (2.3), le foncteur $H^0(A,.)$ est exact sur la catégorie des $i_s^{-1}(\cO_X)$-Modules (T.F. II 3.7.1, 3.4.2, 4.4.3), donc l'homomorphisme
\[
  H^0\bigl(A,i_s^{-1}(E)\bigr)
  \longrightarrow
  H^0\bigl(A,i_s^*(E)\bigr)
\]
est un épimorphisme. Comme $i_s^{-1}(E)$ est mou, l'homomorphisme
\[
  H^0\bigl(X_s,i_s^{-1}(E)\bigr)
  \longrightarrow
  H^0\bigl(A,i_s^{-1}(E)\bigr)
\]
est un épimorphisme, et l'on gagne.

\textbf{Corollaire 2.5.} Soient $p:X\to S$ une variété mixte, et $s\in S$. On suppose que $s$ admet un système fondamental de voisinages paracompacts $S'$ tels que les $p^{-1}(S')$ forment un système fondamental de voisinages paracompacts de $X_s$ (condition réalisée par exemple si $S$ est localement compact et $p$ une application propre d'espaces topologiques). Alors, pour tout $\cO_X$-Module $E$, on a
\[
  R^i p_*(E)_s=0 \quad \text{pour } i>0.
\]

